\section{Implementación de MA-AIRL}

\subsection{Objetivo y alcance}
Se implementó una versión funcional del algoritmo Multi-Agent Adversarial Inverse Reinforcement Learning (MA-AIRL) siguiendo el marco descrito en \texttt{src/MAIRL.pdf}. El objetivo es aprender, por cada clúster de conductores, una política estocástica y una función de recompensa que expliquen las demostraciones observadas en los datos agregados de clústeres, y luego exportar parámetros de \textit{vTypes} para SUMO.

\subsection{Arquitectura general}
La implementación se dividió en dos capas:
\begin{itemize}
\item \textbf{Núcleo MA-AIRL} en \texttt{src/marl\_core.py}: datos expertos, normalización, redes (política, recompensa y potencial), entrenamiento y exportación de parámetros.
\item \textbf{Interfaz} en \texttt{src/multi\_agent\_rl\_app.py}: carga de datos, selección de \textit{features}, control de hiperparámetros, ejecución del entrenamiento y visualización de métricas.
\end{itemize}

\subsection{Datos expertos y construcción del estado}
Los datos de entrada son archivos CSV de clústeres con la columna \texttt{cluster\_label}. Cada fila se interpreta como una muestra experta, con un estado $s$ compuesto por columnas numéricas seleccionadas (por defecto: \texttt{avg\_speed\_kmh}, \texttt{avg\_headway\_s}, \texttt{lane\_change\_rate}). Se implementó:
\begin{itemize}
\item \textbf{Selección de features} desde la UI para definir el vector de estado.
\item \textbf{Estandarización} con un \texttt{StandardScaler} propio (media y desviación por columna).
\item \textbf{Separación por agente} agrupando por \texttt{cluster\_label}.
\end{itemize}

\subsection{Acciones expertas (proxy de SUMO)}
MA-AIRL requiere pares $(s, a)$. Como el CSV no contiene acciones reales, se construyó un \textbf{proxy de acción} para cada muestra:
\begin{itemize}
\item Se mapean métricas a parámetros de car-following: \texttt{maxSpeed}, \texttt{accel}, \texttt{decel}, \texttt{minGap}, \texttt{sigma}, \texttt{impatience}.
\item Se aplican transformaciones heurísticas y se recortan a los límites de \texttt{PARAM\_BOUNDS}.
\item Las acciones se normalizan a $[-1,1]$ para entrenamiento estable.
\end{itemize}
Esto permite entrenar MA-AIRL aunque no existan trayectorias con acciones explícitas.

\subsection{Modelos de MA-AIRL}
Para cada agente (clúster) se instancian tres redes:
\begin{itemize}
\item \textbf{Política gaussiana} $q_\theta(a \mid s)$: MLP que produce $\mu$ y $\log\sigma$.
\item \textbf{Estimador de recompensa} $g_\omega(s,a)$: MLP sobre el par estado-acción.
\item \textbf{Potencial} $h_\phi(s)$: MLP sobre el estado.
\end{itemize}
Se implementó la forma estructurada:
\[
f_{\omega,\phi}(s,a,s') = g_\omega(s,a) + \gamma h_\phi(s') - h_\phi(s)
\]

\subsection{Función discriminadora y objetivo}
Siguiendo la Ecuación (11) del paper, el discriminador se define como:
\[
D_{\omega}(s,a)=\frac{\exp(f_{\omega}(s,a))}{\exp(f_{\omega}(s,a)) + q_\theta(a \mid s)}
\]
El entrenamiento usa:
\begin{itemize}
\item \textbf{Paso del discriminador}: maximiza $\mathbb{E}_{X_E}[\log D] + \mathbb{E}_{X_\pi}[\log(1-D)]$.
\item \textbf{Paso del generador (política)}: maximiza $\mathbb{E}_{X_\pi}[f_\omega(s,a) - \log q_\theta(a \mid s)]$ (Ecuación 12).
\end{itemize}

\subsection{Proxy temporal}
Los datos no incluyen transiciones temporales explícitas $(s,s')$. Para mantener la estructura MA-AIRL se aplicó un \textbf{proxy de un paso}:
\[
s' = s
\]
De esta forma se conserva la forma del potencial sin requerir secuencias completas. Esta decisión se documentó en el código como aproximación.

\subsection{Métricas de entrenamiento}
Durante el entrenamiento se reportan:
\begin{itemize}
\item \textbf{disc\_loss}: pérdida del discriminador.
\item \textbf{policy\_loss}: pérdida de la política.
\item \textbf{reward\_mean}: promedio de $f_\omega$ en muestras expertas.
\end{itemize}
Se guardan en memoria para graficado y análisis posterior.

\subsection{Exportación de parámetros a SUMO}
Al finalizar, se obtienen parámetros por agente evaluando la política en el estado medio y desnormalizando a rangos SUMO. Se exporta un archivo:
\begin{itemize}
\item \texttt{mairl\_vtypes.add.xml} con un \texttt{vType} por clúster.
\end{itemize}
Esto permite inyectar los parámetros aprendidos en simulaciones SUMO reales.

\subsection{Integración en la UI}
Se añadió un selector de algoritmo:
\begin{itemize}
\item \textbf{MA-AIRL (paper)}: entrenamiento adversarial con datos expertos.
\item \textbf{CEM (legacy)}: método anterior basado en cross-entropy con SUMO.
\end{itemize}
También se agregó control de hiperparámetros (batch size, $\gamma$, \textit{learning rates}, tamaño de red y dispositivo).

\subsection{Archivos modificados}
\begin{itemize}
\item \texttt{src/marl\_core.py}: núcleo MA-AIRL (datos, redes, entrenamiento, exportación).
\item \texttt{src/multi\_agent\_rl\_app.py}: UI y flujo de entrenamiento MA-AIRL.
\end{itemize}

\subsection{Limitaciones actuales y próximos pasos}
\begin{itemize}
\item La implementación usa $s'=s$ por falta de secuencias temporales explícitas.
\item Las acciones expertas son heurísticas derivadas de métricas agregadas.
\item El siguiente paso recomendado es construir trayectorias reales desde SUMO o datos temporales para estimar $s'$ y acciones reales.
\end{itemize}
